% Pillar-7: limitations + proposed decisive experiments
\section{Limitations and Proposed Decisive Experiments}
\label{sec:p7-limitations}

\subsection{Limitations}
\label{sec:p7-limitations-list}
Every interventional claim in this paper is bounded by four limitations,
stated in load-bearing order.

\emph{No prospective controller advantage.} The long controller sections are
retrospective replay, analytic projection, or proxy optimization. Iter-199
freezes each prompt's step-0 success estimate, so it is not a closed learning
loop. Iter-203 pools rewards from different prompts and therefore cannot emulate
a larger within-prompt GRPO group. Neither analysis estimates held-out policy
improvement. The required promotion test is a seed-paired, fixed-token comparison
against static $G{=}16$ and naive boundary heuristics.

\emph{Single task.} The E3 audit and the E1 gradient probe each run on one
task family (GSM8K-style and synthetic arithmetic respectively). The theory
is task-agnostic, but the controller's measured competitiveness is not: on
harder or drifting prompt populations the trade of
Table~\ref{tab:p7-e3} could tip in any direction.

\emph{Small $n$.} Four arms, one open trainer, shared seeds, held-out
deltas separated by at most $0.075$: the E3 differences between the top
arms are within the noise band the companion papers' power analysis
(Table~\ref{tab:power-analysis}) would demand for a Tier-A claim. We
therefore claim mechanism-consistent telemetry and competitiveness, not
superiority.

\emph{LoRA-only closed stack for one comparison arm.} The closed-stack
telemetry cited for the label-flip caution comes from LoRA
\citep{hu2022lora} fine-tuning on a stack whose sampler and backward pass we
cannot inspect; the open reimplementation is the only arm-for-arm auditable
setting, and cross-stack numbers are quoted only as motivation, never as
controller evidence.

\emph{Toy-scale gradient evidence.} T3's empirical support is the
directional E1 result ($+0.71$ on a 0.5B model); the coefficient is not
expected to transfer, only the sign and ordering.

We also flag a provenance boundary and two scope choices. The 368-run
audit, the model~$\times$~$G$ grid, the 12-cell head-to-head, and the
E1/E3 runs are internal program records (W\&B projects \texttt{zvf-audit}
and \texttt{zvf-colab-experiments}, June 2026): they were run under a
declared protocol and are summarized faithfully, but they have not passed
through the released-artifact pipeline that covers the repo-internal
sweeps and reward-tensor analyses (Appendix~\ref{sec:p7-appendix},
provenance note); the companion Pillar-5 paper flags the same records the
same way. The scope choices: PCD is validated as a control-loop input
(invariance, shape, aliasing resolution), \emph{not} as an outcome
predictor; and no claim in this paper depends on the veRL/OpenRLHF rows of
the framework comparison, which remain dry-run placeholders
(\tableref{tab:implementations}).

\subsection{Proposed decisive experiments}
\label{sec:p7-proposed}
The following experiments would each settle a claim this paper leaves
directional. They extend the frontier backlog maintained in the repository
(\texttt{platform_hybrid/experiments/FRONTIER\_EXPERIMENT\_BACKLOG.md}).

\textbf{D1: The PCD horse race.} Log per-group reward tensors on
\emph{every} anchor run (not only GSM8K), then test early-window PCD
against held-out outcome at the pre-registered $\rho \geq 0.45$ bar, with
mean reward as the mandatory baseline covariate. This is the experiment the
$\rho = 0.95$ vs $0.56$ proxy result gestures at but cannot replace.

\textbf{D2: Adaptive-$G$ on a hard population.} Rerun the four-arm E3
design on a prompt population with baseline $p \in [0.05, 0.3]$ (e.g.\
competition math), $n \geq 5$ seeds per arm, fixed total rollout budget
rather than fixed steps. Theory predicts this is where fixed-$G$ starves
and the controller's targeted spend should separate from both plain GRPO
and always-on dynamic sampling; a null here would bound the controller's
value to telemetry only.

\textbf{D3: T3 at production scale.} Per-group gradient attribution on an
8B model with an open backward pass, 3+ seeds, reporting
$\mathrm{corr}(\|\nabla\|, \hat p(1-\hat p))$ with confidence intervals and
an ablation over advantage normalization (plain, Dr.~GRPO-style, GSPO-style
\citep{qwen2025gspo}). This either promotes T3 from directional to
quantitative or localizes where the binomial model breaks.

\textbf{D4: Escalate-vs-resample crossover.} A two-dimensional sweep
(prompt difficulty $\times$ intervention aggressiveness) mapping where
adaptive-$G$, dynamic sampling \citep{yu2025dapo}, selective rollouts
\citep{zhang2026greso}, and zero-variance advantage shaping
\citep{le2025rlzvp} each win per marginal rollout---turning
Rule~5 of Section~\ref{sec:p7-rules} from a caution into a decision
boundary.

\textbf{D5: Jitter red-team.} Apply the micro-jitter probe to every
pipeline in the registry: any stack whose reported ZVF changes by more than
$0.05$ under $\eps \sim \mathrm{U}(0, 10^{-4})$ reward jitter is flagged as
reporting tie-artifact telemetry. This converts the falsification of
Section~\ref{sec:p7-jitter} into a standing audit.
